\documentclass[11pt,a4paper]{article}

% ── Packages ──────────────────────────────────────────────────────────────
\usepackage[utf8]{inputenc}
\usepackage[T1]{fontenc}
\usepackage{amsmath,amssymb,amsthm}
\usepackage{mathtools}
\usepackage{bm}
\usepackage{geometry}
\geometry{margin=1in}
\usepackage{hyperref}
\hypersetup{colorlinks=true,linkcolor=blue,citecolor=blue,urlcolor=blue}
\usepackage{listings}
\usepackage{xcolor}
\usepackage{titlesec}
\usepackage{enumitem}
\usepackage{booktabs}
\usepackage{fancyhdr}
\usepackage{graphicx}

% ── Code listing style ───────────────────────────────────────────────────
\definecolor{codebg}{RGB}{248,248,248}
\definecolor{codeframe}{RGB}{200,200,200}
\definecolor{codecomment}{RGB}{0,128,0}
\definecolor{codestring}{RGB}{163,21,21}
\definecolor{codekeyword}{RGB}{0,0,200}

\lstdefinestyle{pythonstyle}{
    language=Python,
    backgroundcolor=\color{codebg},
    frame=single,
    rulecolor=\color{codeframe},
    basicstyle=\ttfamily\small,
    keywordstyle=\color{codekeyword}\bfseries,
    commentstyle=\color{codecomment}\itshape,
    stringstyle=\color{codestring},
    showstringspaces=false,
    breaklines=true,
    breakatwhitespace=false,
    tabsize=4,
    captionpos=b,
    numbers=left,
    numberstyle=\tiny\color{gray},
    numbersep=8pt,
    xleftmargin=1.5em,
    framexleftmargin=1.5em,
    aboveskip=1em,
    belowskip=1em,
    morekeywords={self,True,False,None,as,with,lambda},
    literate={→}{$\rightarrow$}1 {←}{$\leftarrow$}1,
}
\lstset{style=pythonstyle}

% ── Header / Footer ──────────────────────────────────────────────────────
\pagestyle{fancy}
\fancyhf{}
\rhead{Cahn--Hilliard PINN --- PyTorch Implementation}
\lhead{Optimizing the Optimizer}
\cfoot{\thepage}

% ── Title ─────────────────────────────────────────────────────────────────
\title{%
  \textbf{Line-by-Line Analysis of the Cahn--Hilliard PINN}\\[4pt]
  \large PyTorch Implementation with YAML-Driven Configuration\\[4pt]
  \normalsize Based on \emph{Optimizing the Optimizer for PINNs}
  (Kiyani et al., Phys.\ Rev.\ E \textbf{106}, 065303, 2022)
}
\author{}
\date{\today}

% ══════════════════════════════════════════════════════════════════════════
\begin{document}
\maketitle
\tableofcontents
\newpage


% ======================================================================
\section{Overview}
\label{sec:overview}
% ======================================================================

This document provides a complete, cell-by-cell explanation of the
\texttt{Cahn\_Hilliard\_Pytorch.ipynb} notebook.  The notebook solves
the two-dimensional \textbf{Cahn--Hilliard equation} (Model~B dynamics)
with a \textbf{Physics-Informed Neural Network (PINN)} and a
\textbf{YAML-driven configuration system} that allows rapid
experimentation with different optimizers, learning rates, warmup
durations, and other training hyper-parameters --- without modifying a
single line of code.

\subsection*{The PDE}

The Cahn--Hilliard equation describes spinodal decomposition and
phase separation in binary mixtures.  Following Equation~(6) of the
reference paper:
\begin{equation}
\boxed{
  \frac{\partial U}{\partial t}
  = D \;\nabla^{2}\!\Big(
      \nabla^{2} U + a_{2}\,U + a_{4}\,U^{3}
    \Big)
}
\label{eq:ch}
\end{equation}
where:
\begin{itemize}[nosep]
  \item $U(t,x,y)$ is the scalar order parameter (phase field),
  \item $D$ is the diffusion/mobility constant,
  \item $a_{2}$ is the quadratic coefficient in the chemical potential,
  \item $a_{4}$ is the quartic coefficient in the chemical potential,
  \item $\nabla^{2} = \partial_{xx} + \partial_{yy}$ is the 2-D Laplacian.
\end{itemize}

The term inside the outer Laplacian is the \emph{chemical potential}
$\mu = \nabla^{2}U + a_{2}\,U + a_{4}\,U^{3}$.  The PDE therefore
contains \textbf{fourth-order spatial derivatives} (a biharmonic
operator $\nabla^{4}U$), which makes it significantly more
challenging for PINNs than second-order PDEs like Burgers' equation.

\subsection*{Training strategy}

The notebook supports a two-phase training pipeline, configurable
entirely from YAML:

\begin{enumerate}
  \item \textbf{Phase~1 --- Adam} (first-order): fast warmup using
        gradient-only updates with learning-rate decay and residual-based
        adaptive sampling (RAD).
  \item \textbf{Phase~2 --- Quasi-Newton} (second-order): one of several
        BFGS variants (SSBroyden2, SSBFGS\_OL, L-BFGS-B, etc.) from
        a patched \texttt{scipy.optimize.minimize}, with warm-started
        inverse Hessian.
\end{enumerate}

The optimizer, warmup length, power transform, batch size, and all
other training parameters are read from a single YAML configuration
file.  A suite of 17~experiment configs is provided for systematic
benchmarking.


% ======================================================================
\section{Cell 1: Imports and Configuration Path}
\label{sec:imports}
% ======================================================================

\subsection*{Purpose}

Load all required libraries and set the single constant
\texttt{CONFIG\_PATH}, which points to the YAML file that controls
every parameter of the experiment.  \emph{This is the only line that
changes between experiments.}

\subsection*{Libraries used}

\begin{center}
\renewcommand{\arraystretch}{1.3}
\begin{tabular}{@{}ll@{}}
\toprule
\textbf{Library} & \textbf{Role} \\
\midrule
\texttt{torch}, \texttt{torch.nn}  & Neural network, autograd \\
\texttt{numpy}                      & Numerical arrays \\
\texttt{scipy.optimize.minimize}    & Quasi-Newton optimizers \\
\texttt{scipy.linalg.cholesky}      & Positive-definiteness check \\
\texttt{scipy.ndimage.gaussian\_filter} & IC smoothing \\
\texttt{scipy.interpolate.RegularGridInterpolator} & IC evaluation \\
\texttt{yaml}                       & Configuration parsing \\
\texttt{matplotlib}                 & Visualization \\
\texttt{time.perf\_counter}         & Wall-clock timing \\
\bottomrule
\end{tabular}
\end{center}

\subsection*{Code}

\begin{lstlisting}
import os                          # filesystem operations (mkdir, path join)
import math                        # sqrt, log, etc. for weight init
import yaml                        # parse YAML configuration files

import numpy as np                 # numerical array operations
import torch                       # deep learning framework
import torch.nn as nn              # neural network modules
import matplotlib.pyplot as plt    # plotting and figure generation
from time import perf_counter      # high-resolution wall-clock timer

# SciPy: quasi-Newton optimization with patched BFGS variants
from scipy.optimize import minimize          # main optimizer entry point
from scipy.linalg import cholesky, LinAlgError  # Hessian PD check

# SciPy: initial condition generation
from scipy.ndimage import gaussian_filter    # smooth the random IC field
from scipy.interpolate import RegularGridInterpolator  # interpolate IC

# PyTorch utility for flat <-> structured parameter conversion
from torch.nn.utils import parameters_to_vector, vector_to_parameters

# ---- CONFIGURATION PATH ---------------------------------------------------
# Change ONLY this line to switch between experiments.
# All tunable parameters (optimizer, LR, warmup, etc.) are read from YAML.
CONFIG_PATH = "configs/cahn_hilliard_canonical.yaml"  # path to YAML config
\end{lstlisting}


% ======================================================================
\section{Cell 2: Load Configuration and Set All Parameters}
\label{sec:config}
% ======================================================================

\subsection*{Purpose}

Parse the YAML configuration file and bind \emph{every} tunable
quantity to a Python variable.  The YAML schema has these sections:

\begin{center}
\renewcommand{\arraystretch}{1.3}
\begin{tabular}{@{}llp{7cm}@{}}
\toprule
\textbf{Section} & \textbf{Key variables} & \textbf{Description} \\
\midrule
\texttt{equation}   & $D$, $a_2$, $a_4$             & PDE coefficients \\
\texttt{domain}     & $x_{\min},x_{\max},y_{\min},y_{\max},t_{\min},t_{\max}$ & Physical domain bounds \\
\texttt{initial\_condition} & type, seed, $\sigma$   & IC generation parameters \\
\texttt{boundary\_conditions} & type, deriv.\ order  & Periodic BC settings \\
\texttt{network}    & layer dims, activation         & MLP architecture \\
\texttt{training.adam}  & epochs, lr, betas, decay   & Adam warmup config \\
\texttt{training.bfgs}  & epochs, method, variant, batch, power & Quasi-Newton config \\
\texttt{sampling}   & $N_{\mathrm{int}}, N_0, N_b$, RAD & Collocation point counts \\
\texttt{loss\_weights} & $\lambda_{\mathrm{pde}}$, $\lambda_{\mathrm{ic}}$, $\lambda_{\mathrm{bc}}$ & Loss term weights \\
\bottomrule
\end{tabular}
\end{center}

\subsection*{Canonical parameter values (Table~I of the paper)}

\begin{center}
\begin{tabular}{lrl}
\toprule
\textbf{Parameter} & \textbf{Value} & \textbf{Meaning} \\
\midrule
$D$, $a_2$, $a_4$     & 1.0 each  & All equation constants unity \\
Domain                 & $128\times 128$, $t\in[0,20]$ & Spatial grid, time horizon \\
IC                     & $U_0 \sim \mathcal{U}(-1,1)$  & Per-pixel uniform random \\
$\sigma_{\mathrm{smooth}}$ & 2.0 & Gaussian smoothing of IC \\
BCs                    & Periodic  & All boundaries wrap \\
Network                & $[3,40,40,40,40,1]$, tanh & $5{,}121$ parameters \\
Adam                   & 2\,000 epochs, lr$=10^{-3}$ & First-order warmup \\
BFGS                   & 10\,000 iters, SSBroyden2    & Second-order convergence \\
$N_{\mathrm{int}}$     & 10\,000    & Interior collocation points \\
$N_0$                  & 1\,000     & IC enforcement points \\
$N_b$                  & 500        & BC pairs per dimension \\
$\lambda_{\mathrm{pde}},\lambda_{\mathrm{ic}},\lambda_{\mathrm{bc}}$ & 1, 5, 5 & Loss weights \\
\bottomrule
\end{tabular}
\end{center}

\subsection*{Code}

\begin{lstlisting}
with open(CONFIG_PATH, "r") as f:  # open the YAML config file for reading
    cfg = yaml.safe_load(f)        # parse YAML into a nested Python dict

# --- Equation coefficients (Eq. 6 in the paper) ---
D  = cfg["equation"]["D"]         # diffusion / mobility constant
a2 = cfg["equation"]["a2"]        # quadratic coeff in chemical potential
a4 = cfg["equation"]["a4"]        # quartic coeff in chemical potential

# --- Spatial / temporal domain ---
x_min = cfg["domain"]["x_min"]    # left boundary of x-domain
x_max = cfg["domain"]["x_max"]    # right boundary of x-domain
y_min = cfg["domain"]["y_min"]    # bottom boundary of y-domain
y_max = cfg["domain"]["y_max"]    # top boundary of y-domain
t_min = cfg["domain"]["t_min"]    # initial time
t_max = cfg["domain"]["t_max"]    # final time
Lx = x_max - x_min               # physical length in the x-direction
Ly = y_max - y_min               # physical length in the y-direction
normalize_inputs = cfg["domain"].get("normalize_inputs", True)
                                  # whether to map inputs to [-1,1]

# --- Initial condition ---
ic_cfg        = cfg["initial_condition"]   # sub-dict for IC settings
ic_type       = ic_cfg["type"]             # IC type ("uniform_random")
ic_seed       = ic_cfg.get("seed", 42)     # RNG seed for reproducibility
ic_grid_nx    = ic_cfg.get("grid_nx", 128) # IC grid resolution in x
ic_grid_ny    = ic_cfg.get("grid_ny", 128) # IC grid resolution in y
ic_smooth_sig = ic_cfg.get("smoothing_sigma", 2.0)
                                           # Gaussian sigma for IC smoothing

# --- Boundary conditions ---
bc_cfg         = cfg["boundary_conditions"]          # BC sub-dict
bc_type        = bc_cfg["type"]                      # "periodic"
bc_enforcement = bc_cfg.get("enforcement", "soft")   # "soft" penalty
bc_deriv_order = bc_cfg.get("derivative_order", 0)   # match derivatives too?

# --- Network architecture ---
net_cfg         = cfg["network"]                     # network sub-dict
layer_dims      = tuple(net_cfg["layer_dims"])       # e.g. (3,40,40,40,40,1)
activation_name = net_cfg.get("activation", "tanh")  # activation function name
output_scaling  = net_cfg.get("output_scaling", 1.0) # multiply final output

# --- Training schedule: Adam ---
adam_cfg       = cfg["training"]["adam"]              # Adam sub-dict
Nepochs_ADAM   = adam_cfg["epochs"]                   # number of Adam iters
adam_lr        = adam_cfg["lr"]                        # initial learning rate
adam_betas     = tuple(adam_cfg["betas"])              # momentum params (b1,b2)
adam_eps       = adam_cfg["eps"]                       # Adam epsilon
lr_decay_rate  = adam_cfg.get("lr_decay_rate", 0.98)  # LR multiplicative decay
lr_decay_steps = adam_cfg.get("lr_decay_steps", 1000) # decay period

# --- Training schedule: BFGS ---
bfgs_cfg        = cfg["training"]["bfgs"]             # BFGS sub-dict
Nbfgs           = bfgs_cfg["epochs"]                  # total BFGS iterations
Nchange         = bfgs_cfg["batch_size"]              # resample every N iters
bfgs_method     = bfgs_cfg["method"]                  # "BFGS" or "L-BFGS-B"
bfgs_variant    = bfgs_cfg.get("variant", "none")     # e.g. "SSBroyden2"
bfgs_init_scale = bfgs_cfg.get("initial_scale", False)
                                                      # Nocedal-Wright H0 scaling
bfgs_power      = bfgs_cfg.get("power", 1.0)         # power transform exponent
bfgs_maxcor     = bfgs_cfg.get("maxcor", 20)         # L-BFGS-B corrections

# --- Sampling ---
samp_cfg   = cfg["sampling"]                          # sampling sub-dict
Nint       = samp_cfg["n_interior"]                   # PDE collocation pts
N0         = samp_cfg["n_initial"]                    # IC enforcement pts
Nb         = samp_cfg["n_boundary"]                   # BC point pairs
Nresample  = samp_cfg["resample_every"]               # resample frequency
rad_cfg    = samp_cfg.get("rad", {})                  # RAD sub-dict
rad_on     = rad_cfg.get("enabled", True)             # toggle RAD on/off
k1         = rad_cfg.get("k1", 1.0)                  # RAD weighting exponent
k2         = rad_cfg.get("k2", 1.0)                  # RAD baseline probability
Nsampling  = rad_cfg.get("n_candidates", 50000)       # RAD candidate pool size
rad_args   = (k1, k2)                                 # pack RAD params as tuple

# --- Loss weights ---
lw      = cfg["loss_weights"]                         # loss weights sub-dict
lam_pde = lw.get("pde", 1.0)                         # weight for PDE residual
lam_ic  = lw.get("ic", 5.0)                          # weight for IC term
lam_bc  = lw.get("bc", 5.0)                          # weight for BC term

# --- Logging ---
log_cfg     = cfg.get("logging", {})                  # logging sub-dict
Nprint      = log_cfg.get("print_every", 100)         # print frequency
RESULTS_DIR = log_cfg.get("results_dir",
                           "results_cahn_hilliard")   # output directory

# --- Global settings ---
SEED      = cfg.get("seed", 2)                        # global random seed
dtype_str = cfg.get("dtype", "float64")               # float precision

# ---------- Apply settings ----------
torch_dtype = torch.float64 if dtype_str == "float64" \
              else torch.float32                       # resolve torch dtype
torch.set_default_dtype(torch_dtype)                   # set global default
torch.manual_seed(SEED)                                # seed torch RNG
np.random.seed(SEED)                                   # seed numpy RNG
DEVICE = torch.device(                                 # auto-detect GPU
    "cuda" if torch.cuda.is_available() else "cpu")
os.makedirs(RESULTS_DIR, exist_ok=True)                # create output dir
\end{lstlisting}


% ======================================================================
\section{Cell 3: Neural Network Definition}
\label{sec:network}
% ======================================================================

\subsection*{Mathematical description}

The network is a fully-connected feed-forward MLP that maps
$(t,x,y) \in \mathbb{R}^3$ to $\hat{U}(t,x,y) \in \mathbb{R}$:
\begin{equation}
  \hat{U}(t,x,y;\bm{\theta})
  = s \cdot \bm{W}^{[5]}\,\sigma\!\bigl(
      \bm{W}^{[4]}\,\sigma\!\bigl(
        \bm{W}^{[3]}\,\sigma\!\bigl(
          \bm{W}^{[2]}\,\sigma\!\bigl(
            \bm{W}^{[1]}\,\mathcal{N}(t,x,y)
            + \bm{b}^{[1]}
          \bigr) + \bm{b}^{[2]}
        \bigr) + \bm{b}^{[3]}
      \bigr) + \bm{b}^{[4]}
    \bigr) + \bm{b}^{[5]}
\label{eq:mlp}
\end{equation}
where:
\begin{itemize}[nosep]
  \item $\mathcal{N}(t,x,y)$ is the \textbf{input normalisation}
        layer, mapping each coordinate from $[\text{lo},\text{hi}]$
        to $[-1,1]$:
        \[
          \mathcal{N}_j(z_j) = \frac{z_j - \tfrac{1}{2}(\text{lo}_j + \text{hi}_j)}
                                    {\tfrac{1}{2}(\text{hi}_j - \text{lo}_j)}
        \]
        For the canonical config:
        $t \in [0,20] \to [-1,1]$,
        $x \in [0,128] \to [-1,1]$,
        $y \in [0,128] \to [-1,1]$.

  \item $\sigma(z) = \tanh(z)$ is the hidden-layer activation.
        Tanh is $C^{\infty}$, which is essential because we
        differentiate through the network up to \textbf{fourth order}.

  \item $s = \texttt{output\_scaling}$ multiplies the final output
        (identity when $s=1$).

  \item $\bm{\theta} = \{(\bm{W}^{[\ell]},\bm{b}^{[\ell]})\}_{\ell=1}^{5}$
        collects all trainable parameters.
\end{itemize}

\subsubsection*{Parameter count}

For \texttt{layer\_dims} = $(3, 40, 40, 40, 40, 1)$:
\[
  |\bm\theta| = (3\!\times\!40 {+} 40) + 3 \times (40\!\times\!40 {+} 40) + (40\!\times\!1 {+} 1)
  = 160 + 4\,920 + 41 = 5{,}121
\]

The dense inverse Hessian $H_k \in \mathbb{R}^{5121 \times 5121}$
requires $5121^2 \times 8 \approx 210$~MB (float64).

\subsubsection*{Weight initialisation}

The final linear layer uses a \emph{variance-scaling} initialiser
(TF-style \texttt{fan\_avg} uniform):
\begin{equation}
  n = \frac{\text{fan\_in} + \text{fan\_out}}{2}, \qquad
  \ell = \sqrt{\frac{3\,s}{n}}, \qquad
  W_{ij} \sim \mathcal{U}(-\ell,\,\ell), \qquad
  b_j = 0
\label{eq:init}
\end{equation}
where $s = 1/\texttt{output\_scaling}^2$.

\subsection*{Code}

\begin{lstlisting}
# ---- Activation function lookup table ------------------------------------
ACTIVATIONS = {                        # map string names to nn.Module classes
    "tanh": nn.Tanh, "relu": nn.ReLU,
    "gelu": nn.GELU, "silu": nn.SiLU,
    "sigmoid": nn.Sigmoid,
}


class InputNormalization(nn.Module):
    """Non-trainable affine layer: maps [lo, hi] -> [-1, 1] per feature."""
    def __init__(self, lo, hi):
        super().__init__()             # initialise nn.Module base class
        lo = torch.tensor(lo, dtype=torch.get_default_dtype())  # lower bounds
        hi = torch.tensor(hi, dtype=torch.get_default_dtype())  # upper bounds
        self.register_buffer("shift", 0.5 * (lo + hi))  # midpoint of range
        self.register_buffer("scale", 2.0 / (hi - lo))  # 1 / half-range

    def forward(self, x):
        return (x - self.shift) * self.scale  # affine map to [-1, 1]


class OutputScaling(nn.Module):
    """Multiply network output by a fixed constant."""
    def __init__(self, s):
        super().__init__()             # initialise nn.Module base class
        self.s = s                     # store scaling factor

    def forward(self, x):
        return x * self.s              # scale the output


def variance_scaling_init(linear, scale=1.0):
    """TF-style VarianceScaling(mode='fan_avg', distribution='uniform')."""
    fi, fo = linear.in_features, linear.out_features  # fan-in, fan-out
    limit = math.sqrt(3.0 * scale / (0.5 * (fi + fo)))  # compute limit
    with torch.no_grad():                    # no gradient tracking for init
        linear.weight.uniform_(-limit, limit)  # uniform init in [-l, l]
        if linear.bias is not None:
            linear.bias.zero_()              # zero-initialise biases


class CahnHilliardNet(nn.Module):
    """
    Configurable MLP for the Cahn-Hilliard PINN.
    Input  : (t, x, y) -- 3 features
    Output : U(t, x, y) -- 1 feature
    """
    def __init__(self, layer_dims, activation_name="tanh",
                 output_scale=1.0, normalize=False,
                 input_bounds=None):
        super().__init__()             # initialise nn.Module base class
        layers = []                    # accumulate layers in a list

        # Optional input normalization
        if normalize and input_bounds is not None:
            lo, hi = input_bounds      # unpack lower and upper bounds
            layers.append(InputNormalization(lo, hi))  # prepend norm layer

        # Build hidden + output layers from layer_dims
        act_cls = ACTIVATIONS.get(activation_name, nn.Tanh)  # resolve class
        for i in range(len(layer_dims) - 1):                 # iterate pairs
            layers.append(                   # add linear transformation
                nn.Linear(layer_dims[i], layer_dims[i + 1]))
            if i < len(layer_dims) - 2:      # hidden layers get activation
                layers.append(act_cls())     # add tanh (or other) activation

        # Output scaling (identity if scale == 1)
        if output_scale != 1.0:
            layers.append(OutputScaling(output_scale))  # append scaling layer

        self.net = nn.Sequential(*layers)    # wrap everything in Sequential

        # Custom init for the LAST linear layer (variance-scaling)
        for m in reversed(list(self.net.modules())):  # walk modules backwards
            if isinstance(m, nn.Linear):     # find the last Linear layer
                sc = 1.0 / (output_scale ** 2) if output_scale != 0 else 1.0
                variance_scaling_init(m, scale=sc)  # apply custom init
                break                        # only initialise the last one

    def forward(self, x):
        return self.net(x)             # forward pass through all layers


# ---- Build the network from config ----------------------------------------
input_bounds = ([t_min, x_min, y_min],
                [t_max, x_max, y_max])    # bounds for input normalisation
net = CahnHilliardNet(                    # instantiate the network
    layer_dims,                           # layer dimensions from YAML
    activation_name=activation_name,      # activation function name
    output_scale=output_scaling,          # output multiplier
    normalize=normalize_inputs,           # toggle input normalisation
    input_bounds=input_bounds,            # physical coordinate bounds
).to(DEVICE)                              # move to GPU if available

n_params = sum(p.numel()                  # count trainable parameters
               for p in net.parameters()
               if p.requires_grad)
hess_gb = n_params ** 2 * 8 / 1e9        # dense Hessian memory in GB
\end{lstlisting}


% ======================================================================
\section{Cell 4: Initial Condition and Sampling Functions}
\label{sec:sampling}
% ======================================================================

\subsection*{Mathematical description}

\subsubsection*{Initial condition}

The paper specifies (Table~I):
\begin{quote}
  ``uniform random distribution with values between $-1$ and $1$''
\end{quote}
on a $128 \times 128$ grid.  Raw per-pixel noise is not smooth, and a
neural network cannot represent it exactly.  We therefore apply a
Gaussian filter with standard deviation $\sigma = 2$ grid cells:
\begin{equation}
  U_0(x,y) = (G_{\sigma} * \tilde{U}_0)(x,y),
  \qquad
  \tilde{U}_{0,ij} \sim \mathcal{U}(-1, 1),
  \qquad
  G_{\sigma}(r) = \frac{1}{2\pi\sigma^2}\,e^{-r^2/(2\sigma^2)}
\label{eq:ic}
\end{equation}
where $*$ denotes 2-D convolution and the filter uses \texttt{mode='wrap'}
to respect periodic boundaries.

The smoothed field is wrapped into a
\texttt{RegularGridInterpolator} with cubic splines and periodic
padding, so that $U_0(x,y)$ can be evaluated at arbitrary $(x,y)$
during training.

\subsubsection*{Collocation point sampling}

Three sets of points are sampled randomly at each training stage:

\begin{enumerate}[nosep]
  \item \textbf{Interior points}
        $\{(t_i,x_i,y_i)\}_{i=1}^{N_{\mathrm{int}}}$,
        $(t_i,x_i,y_i) \sim \mathcal{U}\bigl([0,20]\times[0,128]\times[0,128]\bigr)$,
        where the PDE residual is enforced.
  \item \textbf{Initial-condition points}
        $\{(0, x_j, y_j)\}_{j=1}^{N_0}$,
        where $u(0,x,y) = U_0(x,y)$ is enforced.
  \item \textbf{Boundary pairs} for periodic BCs:
        $(t,x_{\min},y) \leftrightarrow (t,x_{\max},y)$ and
        $(t,x,y_{\min}) \leftrightarrow (t,x,y_{\max})$, $N_b/2$
        pairs per dimension.
\end{enumerate}

\subsubsection*{Residual-Adaptive Distribution (RAD)}

Collocation points are resampled periodically, weighted by the PDE
residual magnitude:
\begin{equation}
  p(t_i,x_i,y_i) \;\propto\;
  \frac{|f(t_i,x_i,y_i)|^{k_1}}{\overline{|f|^{k_1}}} + k_2
\label{eq:rad}
\end{equation}
This concentrates points in regions of high residual --- where the
network is performing worst --- speeding convergence.

\subsection*{Code}

\begin{lstlisting}
def build_ic_interpolator(ic_cfg, x_min, x_max, y_min, y_max):
    """Generate a smoothed random IC field and return an interpolator."""
    rng   = np.random.RandomState(             # create seeded RNG
                ic_cfg.get("seed", 42))
    nx    = ic_cfg.get("grid_nx", 128)         # grid resolution in x
    ny    = ic_cfg.get("grid_ny", 128)         # grid resolution in y
    sigma = ic_cfg.get("smoothing_sigma", 2.0) # Gaussian sigma
    low   = ic_cfg.get("low", -1.0)            # lower bound for U(-1,1)
    high  = ic_cfg.get("high", 1.0)            # upper bound for U(-1,1)

    field = rng.uniform(low, high,             # raw random IC on grid
                        size=(ny, nx))

    if sigma > 0:                              # apply Gaussian smoothing
        field = gaussian_filter(field,         # smooth the field
                    sigma=sigma, mode="wrap")  # wrap = periodic boundaries

    # Grid coordinates (cell centres)
    dx = (x_max - x_min) / nx                 # grid spacing in x
    dy = (y_max - y_min) / ny                 # grid spacing in y
    xs = np.linspace(x_min + 0.5*dx,          # x cell-centre coordinates
                     x_max - 0.5*dx, nx)
    ys = np.linspace(y_min + 0.5*dy,          # y cell-centre coordinates
                     y_max - 0.5*dy, ny)

    # Pad with periodic wrapping for edge interpolation
    xs_ext = np.concatenate([[x_min], xs, [x_max]])  # extended x coords
    ys_ext = np.concatenate([[y_min], ys, [y_max]])  # extended y coords
    field_ext = np.pad(field,                  # pad field with wrap mode
                       ((1,1),(1,1)), mode="wrap")

    interp = RegularGridInterpolator(          # build cubic interpolator
        (ys_ext, xs_ext), field_ext,           # axes: (y, x) ordering
        method="cubic", bounds_error=False,
        fill_value=None)
    return interp, field                       # return interpolator + raw field


# Build the IC interpolator from config
ic_interp, ic_field = build_ic_interpolator(
    ic_cfg, x_min, x_max, y_min, y_max)


def eval_ic(X_t0):
    """Evaluate IC at (x,y) points. Input (N,3), output (N,1)."""
    x_np = X_t0[:, 1].detach().cpu().numpy()   # extract x coordinates
    y_np = X_t0[:, 2].detach().cpu().numpy()   # extract y coordinates
    vals = ic_interp(                           # query the interpolator
               np.column_stack([y_np, x_np]))  # note: interp expects (y,x)
    return torch.tensor(vals,                   # convert back to torch
               dtype=torch.get_default_dtype()
           ).reshape(-1, 1)                     # shape: (N, 1)


def sample_interior(n):
    """Sample random (t, x, y) inside the domain."""
    t = t_min + (t_max - t_min) * np.random.rand(n)  # random t in [0,20]
    x = x_min + Lx * np.random.rand(n)               # random x in [0,128]
    y = y_min + Ly * np.random.rand(n)               # random y in [0,128]
    return torch.tensor(                              # pack into tensor
        np.column_stack([t, x, y]),
        dtype=torch.get_default_dtype())


def sample_initial(n):
    """Sample (t_min, x, y) for IC enforcement."""
    t = np.full(n, t_min)                             # all t = 0
    x = x_min + Lx * np.random.rand(n)               # random x
    y = y_min + Ly * np.random.rand(n)               # random y
    return torch.tensor(
        np.column_stack([t, x, y]),
        dtype=torch.get_default_dtype())


def sample_boundary_periodic(n):
    """Sample paired boundary points for periodic BC enforcement.

    Returns (X_xlo, X_xhi, X_ylo, X_yhi) -- 4 tensors of shape (n//2, 3).
    """
    n2 = max(n // 2, 1)                   # half the boundary points per dim

    # x-periodic pairs: same (t, y), opposing x faces
    t_x = t_min + (t_max-t_min) * np.random.rand(n2)  # random times
    y_x = y_min + Ly * np.random.rand(n2)             # random y coords
    X_xlo = torch.tensor(                  # points at x = x_min
        np.column_stack([t_x, np.full(n2, x_min), y_x]),
        dtype=torch.get_default_dtype())
    X_xhi = torch.tensor(                  # matched points at x = x_max
        np.column_stack([t_x, np.full(n2, x_max), y_x]),
        dtype=torch.get_default_dtype())

    # y-periodic pairs: same (t, x), opposing y faces
    t_y = t_min + (t_max-t_min) * np.random.rand(n2)  # random times
    x_y = x_min + Lx * np.random.rand(n2)             # random x coords
    X_ylo = torch.tensor(                  # points at y = y_min
        np.column_stack([t_y, x_y, np.full(n2, y_min)]),
        dtype=torch.get_default_dtype())
    X_yhi = torch.tensor(                  # matched points at y = y_max
        np.column_stack([t_y, x_y, np.full(n2, y_max)]),
        dtype=torch.get_default_dtype())

    return X_xlo, X_xhi, X_ylo, X_yhi     # return all 4 boundary tensors
\end{lstlisting}


% ======================================================================
\section{Cell 5: PDE Residual --- Fourth-Order Autograd}
\label{sec:residual}
% ======================================================================

\subsection*{Mathematical description}

\subsubsection*{Expanding the PDE}

The Cahn--Hilliard equation (Eq.~\eqref{eq:ch}) can be expanded in 2-D as:
\begin{equation}
  \frac{\partial U}{\partial t}
  = D\,\Bigl[\,
    \underbrace{\nabla^{4} U}_{\text{biharmonic}}
    + a_2\,\underbrace{\nabla^{2} U}_{\text{Laplacian}}
    + a_4\,\underbrace{\nabla^{2}(U^{3})}_{\text{nonlinear term}}
  \,\Bigr]
\label{eq:ch_expanded}
\end{equation}

\paragraph{Biharmonic operator.}
\begin{equation}
  \nabla^{4} U = U_{xxxx} + 2\,U_{xxyy} + U_{yyyy}
\label{eq:biharm}
\end{equation}

\paragraph{Laplacian.}
\begin{equation}
  \nabla^{2} U = U_{xx} + U_{yy}
\label{eq:laplacian}
\end{equation}

\paragraph{Nonlinear term $\nabla^{2}(U^3)$.}
Rather than computing $U^3$ and differentiating it twice (which would
require 3 additional autograd calls), we expand analytically using the
chain rule:
\begin{align}
  \frac{\partial^2}{\partial x^2}(U^3)
  &= 6\,U\,(U_x)^2 + 3\,U^2\,U_{xx} \label{eq:u3xx}\\
  \frac{\partial^2}{\partial y^2}(U^3)
  &= 6\,U\,(U_y)^2 + 3\,U^2\,U_{yy} \label{eq:u3yy}
\end{align}
Therefore:
\begin{equation}
  \nabla^{2}(U^3) = 6\,U\,(U_x^2 + U_y^2) + 3\,U^2\,(U_{xx} + U_{yy})
\label{eq:lap_u3}
\end{equation}

\subsubsection*{Derivative chain via autograd}

Computing all required derivatives takes \textbf{8 calls} to
\texttt{torch.autograd.grad}:

\begin{center}
\renewcommand{\arraystretch}{1.3}
\begin{tabular}{@{}clll@{}}
\toprule
\textbf{Call} & \textbf{Differentiand} & \textbf{Extracts} & \textbf{Order} \\
\midrule
1 & $U$     & $U_t,\; U_x,\; U_y$      & 1st \\
2 & $U_x$   & $U_{xx}$                  & 2nd \\
3 & $U_y$   & $U_{yy}$                  & 2nd \\
4 & $U_{xx}$ & $U_{xxx},\; U_{xxy}$     & 3rd \\
5 & $U_{yy}$ & $U_{yyy}$                & 3rd \\
6 & $U_{xxx}$ & $U_{xxxx}$              & 4th \\
7 & $U_{xxy}$ & $U_{xxyy}$             & 4th \\
8 & $U_{yyy}$ & $U_{yyyy}$              & 4th \\
\bottomrule
\end{tabular}
\end{center}

Every call uses \texttt{create\_graph=True} so that second-order
optimizers can backpropagate through the residual.

\subsubsection*{Final residual}

Combining everything:
\begin{equation}
\boxed{
  f(t,x,y;\bm\theta)
  = U_t - D\,\bigl[\,
    (U_{xxxx} + 2\,U_{xxyy} + U_{yyyy})
    + a_2\,(U_{xx} + U_{yy})
    + a_4\,\bigl(6\,U(U_x^2 + U_y^2) + 3\,U^2(U_{xx} + U_{yy})\bigr)
  \,\bigr]
}
\label{eq:residual}
\end{equation}

The PDE is satisfied when $f \equiv 0$ everywhere.

\subsection*{Code}

\begin{lstlisting}
def forward_pass(net, X):
    """Network forward pass. Returns U of shape (N, 1)."""
    return net(X)[:, 0:1]             # keep (N,1) shape for autograd


def compute_pde_residual(net, X):
    """Compute the Cahn-Hilliard PDE residual at collocation points.

    Args:
        net : the neural network
        X   : tensor (N, 3) of (t, x, y) points
    Returns:
        U        : network prediction (N, 1)
        residual : PDE residual       (N,)
    """
    X = X.clone().detach().requires_grad_(True)  # enable gradient tracking
    U = forward_pass(net, X)                     # forward pass: get U
    ones = torch.ones_like(U)                    # grad_outputs for (N,1)
    ones1 = torch.ones_like(U[:, 0])             # grad_outputs for (N,)

    # ---- 1st derivatives: U_t, U_x, U_y -----------------------------------
    dU  = torch.autograd.grad(U, X,             # differentiate U w.r.t. X
              grad_outputs=ones,                 # vector-Jacobian product
              create_graph=True,                 # keep graph for higher derivs
              retain_graph=True)[0]              # [0] extracts the tensor
    U_t = dU[:, 0]                               # time derivative
    U_x = dU[:, 1]                               # x spatial derivative
    U_y = dU[:, 2]                               # y spatial derivative

    # ---- 2nd derivatives: U_xx, U_yy ---------------------------------------
    dUx = torch.autograd.grad(U_x, X,           # differentiate U_x w.r.t. X
              grad_outputs=ones1,
              create_graph=True,
              retain_graph=True)[0]
    U_xx = dUx[:, 1]                             # d^2U/dx^2

    dUy = torch.autograd.grad(U_y, X,           # differentiate U_y w.r.t. X
              grad_outputs=ones1,
              create_graph=True,
              retain_graph=True)[0]
    U_yy = dUy[:, 2]                             # d^2U/dy^2

    # ---- 3rd derivatives: U_xxx, U_xxy, U_yyy -----------------------------
    dUxx = torch.autograd.grad(U_xx, X,          # differentiate U_xx w.r.t. X
               grad_outputs=ones1,
               create_graph=True,
               retain_graph=True)[0]
    U_xxx = dUxx[:, 1]                            # d^3U/dx^3
    U_xxy = dUxx[:, 2]                            # d^3U/dx^2 dy

    dUyy = torch.autograd.grad(U_yy, X,          # differentiate U_yy w.r.t. X
               grad_outputs=ones1,
               create_graph=True,
               retain_graph=True)[0]
    U_yyy = dUyy[:, 2]                            # d^3U/dy^3

    # ---- 4th derivatives: U_xxxx, U_xxyy, U_yyyy --------------------------
    dUxxx = torch.autograd.grad(U_xxx, X,         # differentiate U_xxx
                grad_outputs=ones1,
                create_graph=True,
                retain_graph=True)[0]
    U_xxxx = dUxxx[:, 1]                           # d^4U/dx^4

    dUxxy = torch.autograd.grad(U_xxy, X,         # differentiate U_xxy
                grad_outputs=ones1,
                create_graph=True,
                retain_graph=True)[0]
    U_xxyy = dUxxy[:, 2]                          # d^4U/dx^2 dy^2

    dUyyy = torch.autograd.grad(U_yyy, X,          # differentiate U_yyy
                grad_outputs=ones1,
                create_graph=True,
                retain_graph=True)[0]
    U_yyyy = dUyyy[:, 2]                           # d^4U/dy^4

    # ---- Assemble PDE terms ------------------------------------------------
    biharmonic = U_xxxx + 2.0*U_xxyy + U_yyyy     # nabla^4 U (Eq. 5)
    laplacian  = U_xx + U_yy                        # nabla^2 U (Eq. 6)

    # nabla^2(U^3) via analytical chain rule (Eq. 7)
    Uf = U[:, 0]                                    # flatten U from (N,1)
    lap_U3 = (6.0 * Uf * (U_x**2 + U_y**2)         # 6U(U_x^2 + U_y^2)
              + 3.0 * Uf**2 * laplacian)             # + 3U^2(U_xx + U_yy)

    # Full RHS of the PDE
    rhs = D * (biharmonic                           # D * nabla^4 U
               + a2 * laplacian                      # + D * a2 * nabla^2 U
               + a4 * lap_U3)                        # + D * a4 * nabla^2(U^3)

    residual = U_t - rhs                             # residual = LHS - RHS
    return U, residual                               # should be ~0 everywhere


def compute_bc_derivatives(net, X):
    """Compute U and first derivatives at boundary points."""
    X = X.clone().detach().requires_grad_(True)  # enable gradient tracking
    U = forward_pass(net, X)                     # forward pass
    ones = torch.ones_like(U)                    # grad_outputs
    dU = torch.autograd.grad(U, X,               # first derivatives
             grad_outputs=ones,
             create_graph=True,
             retain_graph=True)[0]
    return U, dU[:, 1], dU[:, 2]  # U, dU/dx, dU/dy
\end{lstlisting}


% ======================================================================
\section{Cell 6: Loss Function, Gradient Computation, and Utilities}
\label{sec:loss}
% ======================================================================

\subsection*{Mathematical description}

The total PINN loss is:
\begin{equation}
\boxed{
  \mathcal{L}(\bm\theta)
  = \lambda_{\mathrm{pde}}
    \underbrace{
      \frac{1}{N_{\mathrm{int}}} \sum_{i=1}^{N_{\mathrm{int}}}
      \bigl[f(t_i,x_i,y_i)\bigr]^2
    }_{\mathcal{L}_{\mathrm{PDE}}}
  + \lambda_{\mathrm{ic}}
    \underbrace{
      \frac{1}{N_0} \sum_{j=1}^{N_0}
      \bigl[\hat{U}(0,x_j,y_j) - U_0(x_j,y_j)\bigr]^2
    }_{\mathcal{L}_{\mathrm{IC}}}
  + \lambda_{\mathrm{bc}}\,\mathcal{L}_{\mathrm{BC}}
}
\label{eq:loss}
\end{equation}

The boundary-condition loss for \textbf{periodic BCs} enforces:
\begin{align}
  \mathcal{L}_{\mathrm{BC}} &=
    \mathrm{MSE}\bigl(\hat{U}(t,x_{\min},y),\;\hat{U}(t,x_{\max},y)\bigr)
    + \mathrm{MSE}\bigl(\hat{U}(t,x,y_{\min}),\;\hat{U}(t,x,y_{\max})\bigr)
    \notag\\
    &\quad + \mathrm{MSE}\bigl(\hat{U}_x\big|_{x_{\min}},\;
                               \hat{U}_x\big|_{x_{\max}}\bigr)
    + \mathrm{MSE}\bigl(\hat{U}_y\big|_{y_{\min}},\;
                               \hat{U}_y\big|_{y_{\max}}\bigr)
\label{eq:bc_loss}
\end{align}
The derivative-matching terms (second line) are included when
\texttt{bc\_deriv\_order} $ \ge 1$ in the config.

\subsubsection*{SciPy-compatible wrapper and power transform}

For the quasi-Newton phase, we need a flat-vector interface compatible
with \texttt{scipy.optimize.minimize}:
\[
  \texttt{fun}(\bm\theta) \;\to\;
  \bigl(\mathcal{L}_{\mathrm{eff}},\;
        \nabla_{\bm\theta}\mathcal{L}_{\mathrm{eff}}\bigr),
  \qquad \bm\theta \in \mathbb{R}^{5121}
\]

The optional \textbf{power transform} minimises
$\mathcal{L}_{\mathrm{eff}} = \mathcal{L}^{1/p}$ instead of
$\mathcal{L}$.  When $p = 2$:
\begin{equation}
  \mathcal{L}_{\mathrm{eff}} = \sqrt{\mathcal{L}},
  \qquad
  \nabla \mathcal{L}_{\mathrm{eff}}
  = \frac{1}{2\sqrt{\mathcal{L}}}\,\nabla\mathcal{L}
\label{eq:power}
\end{equation}

By Theorem~1 of the paper, this reduces the Hessian condition number:
\begin{equation}
  \kappa\!\left(\nabla^2\sqrt{\mathcal{L}}\right)
  \;\approx\; \sqrt{\kappa\!\left(\nabla^2\mathcal{L}\right)}
\label{eq:kappa}
\end{equation}

\subsection*{Code}

\begin{lstlisting}
mse = nn.MSELoss()                         # mean-squared-error loss function


def compute_loss(net, X_int, X_ic, bc_data):
    """Full PINN loss for the Cahn-Hilliard equation.

    Args:
        net     : the neural network
        X_int   : (N_int, 3) interior collocation points
        X_ic    : (N_ic, 3) initial-condition points
        bc_data : tuple (X_xlo, X_xhi, X_ylo, X_yhi)
    Returns:
        total_loss : scalar tensor (differentiable)
        pde_loss   : scalar float (for logging)
    """
    X_xlo, X_xhi, X_ylo, X_yhi = bc_data  # unpack boundary pairs

    # ---- PDE residual loss -------------------------------------------------
    _, residual = compute_pde_residual(      # get PDE residual at interior pts
                      net, X_int)
    zeros_r = torch.zeros_like(residual)     # target = 0 everywhere
    L_pde = mse(residual, zeros_r)           # MSE of PDE residual

    # ---- Initial-condition loss --------------------------------------------
    U_ic_pred = forward_pass(net, X_ic)      # network prediction at t=0
    U_ic_true = eval_ic(X_ic).to(            # true IC values from interp.
                    U_ic_pred.device)
    L_ic = mse(U_ic_pred, U_ic_true)         # MSE of IC mismatch

    # ---- Periodic BC loss --------------------------------------------------
    if bc_deriv_order >= 1:                  # match values AND derivatives
        U_lo_x, Ux_lo, _ = compute_bc_derivatives(  # left x-boundary
                                net, X_xlo)
        U_hi_x, Ux_hi, _ = compute_bc_derivatives(  # right x-boundary
                                net, X_xhi)
        U_lo_y, _, Uy_lo  = compute_bc_derivatives(  # bottom y-boundary
                                net, X_ylo)
        U_hi_y, _, Uy_hi  = compute_bc_derivatives(  # top y-boundary
                                net, X_yhi)
        L_bc = (mse(U_lo_x, U_hi_x)         # U(x_min) == U(x_max)
                + mse(U_lo_y, U_hi_y)        # U(y_min) == U(y_max)
                + mse(Ux_lo, Ux_hi)          # dU/dx matches at x bounds
                + mse(Uy_lo, Uy_hi))         # dU/dy matches at y bounds
    else:                                    # match values only
        U_lo_x = forward_pass(net, X_xlo)   # U at left x-boundary
        U_hi_x = forward_pass(net, X_xhi)   # U at right x-boundary
        U_lo_y = forward_pass(net, X_ylo)   # U at bottom y-boundary
        U_hi_y = forward_pass(net, X_yhi)   # U at top y-boundary
        L_bc = (mse(U_lo_x, U_hi_x)         # periodic matching in x
                + mse(U_lo_y, U_hi_y))       # periodic matching in y

    # ---- Weighted total loss -----------------------------------------------
    total = (lam_pde * L_pde                 # PDE residual term
             + lam_ic * L_ic                 # initial condition term
             + lam_bc * L_bc)                # boundary condition term
    return total, float(L_pde.detach())      # return loss + PDE component


def adam_step(net, optimizer, X_int, X_ic, bc_data):
    """One Adam training step: forward -> loss -> backward -> update."""
    for p in net.parameters():               # clear all existing gradients
        p.grad = None
    loss_val, pde_val = compute_loss(        # compute the full PINN loss
                            net, X_int, X_ic, bc_data)
    loss_val.backward()                      # backprop to compute gradients
    optimizer.step()                         # Adam parameter update
    return float(loss_val.detach()), pde_val # return scalar loss values


def scipy_loss_and_grad(weights, net, X_int, X_ic,
                        bc_data, power=1.0):
    """Flat numpy interface for scipy.optimize.minimize.

    1. Load weights into network
    2. Compute loss + gradients via autograd
    3. Return (scalar_loss, flat_gradient) as numpy arrays
    """
    device = next(net.parameters()).device   # get the device (CPU or GPU)
    dtype  = next(net.parameters()).dtype    # get the dtype (float64)
    w = torch.as_tensor(weights,            # convert flat numpy -> torch
                        dtype=dtype,
                        device=device)
    with torch.no_grad():                   # load weights into the network
        vector_to_parameters(w,             # assign flat vector to params
                             net.parameters())

    loss_val, _ = compute_loss(             # compute full PINN loss
                      net, X_int, X_ic, bc_data)

    # Optional power transform: minimise L^(1/p) instead of L
    loss_eff = loss_val if power == 1.0 \
               else loss_val ** (1.0 / power)  # apply power transform

    params = list(net.parameters())         # list of parameter tensors
    grads = torch.autograd.grad(            # compute gradient of loss_eff
                loss_eff, params,           #   w.r.t. all network params
                create_graph=False,         #   no higher-order tracking
                retain_graph=False,         #   free the graph after
                allow_unused=False)         #   all params must be used
    g_flat = torch.cat(                     # flatten gradient list
                 [g.reshape(-1) for g in grads]
             ).detach().cpu().numpy()       # convert to numpy for SciPy
    return float(loss_eff.detach().cpu()),  # return (loss, grad) tuple
           g_flat


def adaptive_rad(net, n_int, rad_args, n_cand=50000):
    """Resample interior points weighted by |PDE residual|.

    High-residual regions get more points -> faster convergence.
    """
    X_cand = sample_interior(n_cand)         # generate candidate points
    k1, k2 = rad_args                        # unpack RAD parameters
    Y = torch.abs(                           # absolute PDE residual
            compute_pde_residual(net, X_cand)[1]
        ).detach().reshape(-1)
    w = Y ** k1                              # weight = |residual|^k1
    p = (w / w.mean() + k2)                  # normalise + baseline prob
    p = (p / p.sum()).clamp_min(1e-12)        # ensure valid probabilities
    ids = torch.multinomial(p,               # weighted random sampling
              num_samples=n_int,
              replacement=False)
    return X_cand[ids]                        # return selected points
\end{lstlisting}


% ======================================================================
\section{Cell 7: Phase 1 --- Adam Training (First-Order Warmup)}
\label{sec:adam}
% ======================================================================

\subsection*{Mathematical description}

Phase~1 runs \texttt{Nepochs\_ADAM} iterations of the Adam optimizer.
The learning rate follows an exponential decay schedule:
\begin{equation}
  \eta_t = \eta_0 \cdot r^{t/s}
  \qquad
  \text{(canonical: } \eta_0 = 10^{-3},\; r = 0.98,\; s = 1000\text{)}
\label{eq:lr_schedule}
\end{equation}

Adam updates (Kingma \& Ba, 2015):
\begin{align}
  \bm{m}_t &= \beta_1\,\bm{m}_{t-1} + (1\!-\!\beta_1)\,\bm{g}_t
    & &\text{(first moment)} \\
  \bm{v}_t &= \beta_2\,\bm{v}_{t-1} + (1\!-\!\beta_2)\,\bm{g}_t^2
    & &\text{(second moment)} \\
  \hat{\bm{m}}_t &= \bm{m}_t / (1\!-\!\beta_1^t),\quad
  \hat{\bm{v}}_t = \bm{v}_t / (1\!-\!\beta_2^t)
    & &\text{(bias correction)} \\
  \bm\theta_{t+1} &= \bm\theta_t
    - \eta_t\,\hat{\bm{m}}_t / (\sqrt{\hat{\bm{v}}_t} + \epsilon)
    & &\text{(update)}
\label{eq:adam}
\end{align}

with $\beta_1 = 0.99$, $\beta_2 = 0.999$, $\epsilon = 10^{-20}$.

Every \texttt{Nresample} epochs, collocation points are regenerated
using RAD (Eq.~\eqref{eq:rad}).  If \texttt{Nepochs\_ADAM} $= 0$
(the \texttt{warmup\_0} config), this phase is skipped entirely.

\emph{Purpose:} first-order methods are cheap per-iteration but
converge slowly near a local minimum.  Adam gets the network into a
reasonable basin from which the more powerful quasi-Newton optimizer
can take over.

\subsection*{Code}

\begin{lstlisting}
# ---- Generate initial training data ---------------------------------------
X_int   = sample_interior(Nint)             # N_int interior collocation pts
X_ic    = sample_initial(N0)                # N_0 initial-condition pts
bc_data = sample_boundary_periodic(Nb)      # paired boundary points

adam_losses = []                             # accumulate loss for plotting
adam_t0 = perf_counter()                     # start wall-clock timer

if Nepochs_ADAM > 0:                         # skip if warmup = 0
    # ---- Configure Adam ----------------------------------------------------
    optimizer = torch.optim.Adam(            # create Adam optimizer
        net.parameters(),                    # parameters to optimise
        lr=adam_lr,                           # initial learning rate
        betas=adam_betas,                     # momentum coefficients
        eps=adam_eps)                         # numerical stability epsilon
    scheduler = torch.optim.lr_scheduler.LambdaLR(   # decay scheduler
        optimizer,
        lr_lambda=lambda step:               # multiply LR by this factor
            lr_decay_rate ** (step / lr_decay_steps))

    # ---- Training loop -----------------------------------------------------
    for epoch in range(Nepochs_ADAM):        # loop over Adam epochs
        # Periodically resample collocation points (RAD if enabled)
        if (epoch + 1) % Nresample == 0:     # check if resample is due
            if rad_on:                       # use Residual-Adaptive Dist.
                X_int = adaptive_rad(net, Nint, rad_args, Nsampling)
            else:                            # uniform random resampling
                X_int = sample_interior(Nint)
            X_ic    = sample_initial(N0)     # resample IC points
            bc_data = sample_boundary_periodic(Nb)  # resample BC pairs

        loss_v, pde_v = adam_step(           # one Adam step
            net, optimizer, X_int, X_ic, bc_data)
        scheduler.step()                     # update learning rate
        adam_losses.append(loss_v)            # record loss

        if (epoch + 1) % Nprint == 0:        # periodic console output
            print(f"Adam  epoch {epoch+1:5d}/{Nepochs_ADAM}  "
                  f"loss={loss_v:.4e}  pde={pde_v:.4e}")

adam_time = perf_counter() - adam_t0          # total Adam wall time
if adam_losses:                               # at least one epoch ran
    print(f"\nAdam phase complete: {adam_time:.1f}s  "
          f"final loss={adam_losses[-1]:.4e}")
else:                                         # warmup_0 config
    print("Adam phase skipped (0 epochs).")
\end{lstlisting}


% ======================================================================
\section{Cell 8: Phase 2 --- Quasi-Newton Training (BFGS / L-BFGS-B)}
\label{sec:bfgs}
% ======================================================================

\subsection*{Mathematical description}

\subsubsection*{Why switch from Adam to BFGS?}

Near the solution, the loss landscape is approximately quadratic.
First-order methods converge linearly with rate
$\approx 1 - 2/\kappa(\nabla^2\!\mathcal{L})$, which can be
arbitrarily slow when the Hessian is ill-conditioned.  Quasi-Newton
methods (BFGS) build an approximate inverse Hessian
$H_k \approx [\nabla^2\!\mathcal{L}]^{-1}$ and achieve
\textbf{superlinear} convergence.

\subsubsection*{Standard BFGS update}

Define the step $\bm{s}_k = \bm\theta_{k+1} - \bm\theta_k$ and the
gradient difference $\bm{y}_k = \nabla\mathcal{L}_{k+1} - \nabla\mathcal{L}_k$:
\begin{equation}
  H_{k+1} = \bigl(I - \rho_k\,\bm{s}_k\,\bm{y}_k^\top\bigr)\,
             H_k\,
             \bigl(I - \rho_k\,\bm{y}_k\,\bm{s}_k^\top\bigr)
             + \rho_k\,\bm{s}_k\,\bm{s}_k^\top
\label{eq:bfgs}
\end{equation}
where $\rho_k = 1 / (\bm{y}_k^\top\bm{s}_k)$.

\subsubsection*{Self-Scaled Broyden (SSBroyden2)}

The paper's key contribution is a \emph{self-scaled} variant that
introduces a scaling factor $\tau_k$ and a Broyden-class parameter
$\phi_k$:
\begin{equation}
  H_{k+1} = \frac{1}{\tau_k}\!\left(
    H_k
    - \frac{H_k\,\bm{y}_k\,\bm{y}_k^\top H_k}
           {\bm{y}_k^\top H_k\,\bm{y}_k}
    + \phi_k\,(\bm{y}_k^\top H_k\,\bm{y}_k)\,\bm{v}_k\,\bm{v}_k^\top
  \right)
  + \rho_k\,\bm{s}_k\,\bm{s}_k^\top
\label{eq:ssbroyden}
\end{equation}

The self-scaling factor $\tau_k$ is chosen to bound the condition
number of $H_{k+1}$, which is why SSBroyden2 converges faster than
vanilla BFGS for ill-conditioned PINN loss landscapes.

\subsubsection*{Supported optimizer variants}

\begin{center}
\renewcommand{\arraystretch}{1.3}
\begin{tabular}{@{}lp{9cm}@{}}
\toprule
\textbf{Variant} & \textbf{Description} \\
\midrule
\texttt{BFGS\_scipy}  & Standard BFGS (no self-scaling) \\
\texttt{SSBFGS\_OL}   & Oren--Luenberger self-scaling \\
\texttt{SSBFGS\_AB}   & Al-Baali self-scaling \\
\texttt{SSBroyden1}   & Broyden class $\phi=1$ (DFP direction) \\
\texttt{SSBroyden2}   & Broyden class with $\phi$ optimised (canonical) \\
\texttt{L-BFGS-B}     & Limited-memory BFGS ($O(n \cdot m)$ memory instead of $O(n^2)$) \\
\bottomrule
\end{tabular}
\end{center}

\subsubsection*{Batched loop with warm-starting}

The optimizer is called in a loop: each call runs at most
\texttt{Nchange} iterations, after which we:
\begin{enumerate}[nosep]
  \item Extract the updated weights $\bm\theta$ and the inverse
        Hessian $H_k$ (dense BFGS only).
  \item Symmetrise $H_k$: $H_k \leftarrow \frac{1}{2}(H_k + H_k^\top)$.
  \item Check positive-definiteness via Cholesky.  If $H_k$ has
        degenerated (numerical errors), reset $H_k \leftarrow I$.
  \item Resample collocation points using RAD.
  \item Continue with warm-started $H_k$.
\end{enumerate}

For L-BFGS-B, there is no dense Hessian to preserve; only the last
\texttt{maxcor} $(s,y)$ pairs are stored.

If \texttt{Nbfgs}~$=0$ (Adam-only config), this cell is a no-op.

\subsection*{Code}

\begin{lstlisting}
if Nbfgs > 0:                                # skip if Adam-only mode
    # ---- Setup -------------------------------------------------------------
    initial_weights = parameters_to_vector(   # flatten all params to numpy
        [p.detach() for p in net.parameters()]
    ).cpu().numpy()

    # Dense Hessian only for BFGS variants (not L-BFGS-B)
    use_dense_hessian = (bfgs_method == "BFGS")  # True for all dense variants
    if use_dense_hessian:
        H0 = np.eye(initial_weights.size,    # initial inverse Hessian = I
                     dtype=np.float64)

    cont = 0                                  # global BFGS iteration counter
    n_ckpts = Nbfgs // Nprint                 # number of logging checkpoints
    bfgs_losses = np.zeros(n_ckpts)           # array for recorded losses
    bfgs_pde    = np.zeros(n_ckpts)           # array for PDE losses

    initial_scale = bfgs_init_scale           # Nocedal-Wright scaling toggle
    power = bfgs_power                        # power transform exponent

    # ---- Callback (runs every BFGS iteration) ------------------------------
    def bfgs_callback(*, intermediate_result):
        """Log loss every Nprint iterations."""
        global cont                           # track total iteration count
        cont += 1                             # increment counter
        if cont % Nprint != 0:                # only log at checkpoints
            return
        idx = cont // Nprint - 1              # compute checkpoint index
        if idx < 0 or idx >= n_ckpts:         # bounds check
            return
        loss_val = float(intermediate_result.fun)  # get current loss
        if power != 1.0:                      # undo power transform for log
            loss_val = loss_val ** power
        bfgs_losses[idx] = loss_val           # store in checkpoint array
        optimizer_label = bfgs_variant \
            if use_dense_hessian else bfgs_method  # label for printing
        print(f"{optimizer_label}  iter "     # console output
              f"{cont:5d}/{Nbfgs}  "
              f"loss={loss_val:.4e}")

    # ---- Build optimizer options dict --------------------------------------
    def build_bfgs_options():
        """Construct options dict for scipy.optimize.minimize."""
        opts = {"maxiter": Nchange, "gtol": 0}  # base options
        if bfgs_method == "BFGS":             # dense BFGS with warm-start
            opts["hess_inv0"]     = H0        # pass current inv. Hessian
            opts["method_bfgs"]   = bfgs_variant  # e.g. "SSBroyden2"
            opts["initial_scale"] = initial_scale  # Noc.-Wright toggle
        elif bfgs_method == "L-BFGS-B":       # limited-memory BFGS
            opts["maxcor"] = bfgs_maxcor      # num correction pairs
            opts["ftol"]   = 0                # no function-tol stopping
        return opts                           # return the completed dict

    # ---- Main BFGS loop ----------------------------------------------------
    bfgs_t0 = perf_counter()                  # start wall-clock timer

    while cont < Nbfgs:                       # loop until budget exhausted
        result = minimize(                    # call SciPy optimizer
            scipy_loss_and_grad,              # objective + gradient function
            initial_weights,                  # current flat weight vector
            args=(net, X_int, X_ic,           # extra args passed to function
                  bc_data, power),
            method=bfgs_method,               # "BFGS" or "L-BFGS-B"
            jac=True,                         # function returns (f, grad)
            options=build_bfgs_options(),      # method-specific options
            tol=0,                            # no early stopping
            callback=bfgs_callback,           # logging callback
        )

        # ---- Extract results from this batch -------------------------------
        initial_weights = result.x            # updated flat weight vector

        if use_dense_hessian:                 # dense BFGS variants only
            H0 = result.hess_inv              # extract inverse Hessian
            H0 = 0.5 * (H0 + H0.T)           # symmetrise (numeric safety)
            try:
                cholesky(H0)                  # check positive-definiteness
            except LinAlgError:               # Hessian has degenerated
                print("  [!] H lost positive-"
                      "definiteness -- "
                      "resetting to I")
                H0 = np.eye(                  # reset to identity
                    len(initial_weights),
                    dtype=np.float64)

        # ---- Resample collocation points (RAD) -----------------------------
        if rad_on:                            # residual-adaptive dist.
            X_int = adaptive_rad(net, Nint,
                        rad_args, Nsampling)
        else:                                 # uniform random resampling
            X_int = sample_interior(Nint)
        X_ic    = sample_initial(N0)          # resample IC points
        bc_data = sample_boundary_periodic(Nb)  # resample BC pairs

        initial_scale = False                 # only on first batch (if ever)

    bfgs_time = perf_counter() - bfgs_t0      # total BFGS wall time
    optimizer_label = bfgs_variant \
        if use_dense_hessian else bfgs_method  # label for printing
    print(f"\n{optimizer_label} phase "       # summary output
          f"complete: {bfgs_time:.1f}s")
    print(f"Total training time: "
          f"{adam_time + bfgs_time:.1f}s")

else:                                         # Adam-only mode (Nbfgs == 0)
    bfgs_time = 0.0                           # no BFGS time
    n_ckpts = 0                               # no checkpoints
    bfgs_losses = np.array([])                # empty loss array
    bfgs_pde    = np.array([])                # empty PDE loss array
    print(f"\nAdam-only mode: no BFGS phase.  "
          f"Total time: {adam_time:.1f}s")
\end{lstlisting}


% ======================================================================
\section{Cell 9: Visualisation and Results}
\label{sec:viz}
% ======================================================================

\subsection*{Purpose}

Generate three diagnostic plots and save them to the results directory:

\begin{enumerate}[nosep]
  \item \textbf{Loss curve}: semi-log plot of total loss across both
        Adam and BFGS phases, with a vertical line at the transition.
  \item \textbf{Solution snapshots}: $U(t,x,y)$ at 5 time slices
        equally spaced from $t=0$ to $t=20$, showing the predicted
        phase-separation dynamics.
  \item \textbf{IC comparison}: side-by-side true $U_0(x,y)$ vs.\
        network prediction at $t=0$, validating that the PINN learned
        the initial condition.
\end{enumerate}

\subsection*{Code}

\begin{lstlisting}
os.makedirs(RESULTS_DIR, exist_ok=True)       # ensure output dir exists
optimizer_label = (                            # build label for plot title
    bfgs_variant if bfgs_method == "BFGS"
    else bfgs_method
) if Nbfgs > 0 else "none"

# ---- 1. Loss curve --------------------------------------------------------
fig, ax = plt.subplots(figsize=(10, 4))       # create figure and axes
ax.semilogy(                                   # plot Adam losses (log-y)
    range(1, len(adam_losses) + 1),            # x = epoch number
    adam_losses, label="Adam", alpha=0.8)       # y = loss value

if Nbfgs > 0:                                 # add BFGS losses if present
    bfgs_epochs = (Nepochs_ADAM                # BFGS epoch numbers
                   + np.arange(1, n_ckpts+1) * Nprint)
    valid = bfgs_losses > 0                    # skip unfilled checkpoints
    if valid.any():                            # plot only valid entries
        ax.semilogy(bfgs_epochs[valid],        # x = BFGS epoch numbers
                    bfgs_losses[valid],         # y = logged losses
                    label=optimizer_label,
                    alpha=0.8)
    ax.axvline(Nepochs_ADAM,                   # vertical line at transition
               color="gray", ls="--", lw=0.8,
               label="Adam->BFGS")

ax.set_xlabel("Iteration")                    # x-axis label
ax.set_ylabel("Total Loss")                   # y-axis label
ax.set_title(f"Cahn-Hilliard PINN -- "        # title with experiment name
             f"{cfg['experiment']['name']}")
ax.legend()                                    # show legend
fig.tight_layout()                             # adjust spacing
fig.savefig(os.path.join(RESULTS_DIR,         # save to disk
            "loss_curve.png"), dpi=150)
plt.show()                                     # display in notebook

# ---- 2. Solution snapshots ------------------------------------------------
n_snap = 5                                     # number of time snapshots
snap_times = np.linspace(t_min, t_max, n_snap)  # evenly spaced times
nx_plot, ny_plot = 64, 64                      # plot grid resolution
xs_plot = np.linspace(x_min, x_max, nx_plot)   # x coordinates for plot
ys_plot = np.linspace(y_min, y_max, ny_plot)   # y coordinates for plot
xx, yy = np.meshgrid(xs_plot, ys_plot)         # 2D meshgrid

fig, axes = plt.subplots(1, n_snap,           # one subplot per snapshot
                         figsize=(4*n_snap, 3.5))
net.eval()                                     # set network to eval mode
for i, t_val in enumerate(snap_times):         # loop over time slices
    tt = np.full_like(xx, t_val)               # constant time array
    X_plot = torch.tensor(                     # build input tensor (N, 3)
        np.column_stack([tt.ravel(),           # t column
                         xx.ravel(),           # x column
                         yy.ravel()]),         # y column
        dtype=torch.get_default_dtype(),
        device=DEVICE)
    with torch.no_grad():                      # no gradient needed for plot
        U_plot = forward_pass(net, X_plot      # predict U at grid points
                 ).cpu().numpy().reshape(
                     ny_plot, nx_plot)          # reshape to 2D grid
    im = axes[i].pcolormesh(xs_plot, ys_plot,  # plot as coloured mesh
             U_plot, cmap="RdBu_r",
             shading="auto")
    axes[i].set_title(f"t = {t_val:.1f}")      # label time slice
    axes[i].set_aspect("equal")                # square aspect ratio
    fig.colorbar(im, ax=axes[i],               # add colour bar
                 fraction=0.046)
net.train()                                    # restore training mode
fig.suptitle("Predicted U(t, x, y)",          # figure title
             fontsize=14, y=1.02)
fig.tight_layout()                             # adjust spacing
fig.savefig(os.path.join(RESULTS_DIR,         # save snapshot figure
            "snapshots.png"), dpi=150,
            bbox_inches="tight")
plt.show()                                     # display in notebook

# ---- 3. IC comparison -----------------------------------------------------
fig, axes = plt.subplots(1, 2,                # two subplots: true vs pred
                         figsize=(9, 3.5))

# True IC from interpolator
X_ic_plot = torch.tensor(                      # build t=0 input tensor
    np.column_stack([np.zeros(nx_plot*ny_plot), # t = 0 everywhere
                     xx.ravel(),               # x grid
                     yy.ravel()]),             # y grid
    dtype=torch.get_default_dtype())
U_ic_true_plot = eval_ic(X_ic_plot             # evaluate true IC
    ).numpy().reshape(ny_plot, nx_plot)        # reshape to 2D
im0 = axes[0].pcolormesh(xs_plot, ys_plot,     # plot true IC
          U_ic_true_plot, cmap="RdBu_r",
          shading="auto")
axes[0].set_title("True IC  U_0(x,y)")        # subplot title
axes[0].set_aspect("equal")                   # square aspect
fig.colorbar(im0, ax=axes[0], fraction=0.046) # colour bar

# Predicted IC from network
net.eval()                                     # eval mode (no dropout etc.)
X_ic_dev = X_ic_plot.to(DEVICE)               # move to device
with torch.no_grad():                          # no gradient tracking
    U_ic_pred_plot = forward_pass(             # network prediction at t=0
        net, X_ic_dev
    ).cpu().numpy().reshape(ny_plot, nx_plot)   # reshape to 2D
net.train()                                    # restore training mode
im1 = axes[1].pcolormesh(xs_plot, ys_plot,     # plot predicted IC
          U_ic_pred_plot, cmap="RdBu_r",
          shading="auto")
axes[1].set_title("PINN at t = 0")            # subplot title
axes[1].set_aspect("equal")                   # square aspect
fig.colorbar(im1, ax=axes[1], fraction=0.046) # colour bar

fig.suptitle("Initial Condition Comparison",  # figure title
             fontsize=14, y=1.02)
fig.tight_layout()                             # adjust spacing
fig.savefig(os.path.join(RESULTS_DIR,         # save IC comparison
            "ic_comparison.png"), dpi=150,
            bbox_inches="tight")
plt.show()                                     # display in notebook

# ---- 4. Summary -----------------------------------------------------------
print(f"\nExperiment: "                        # print experiment name
      f"{cfg['experiment']['name']}")
print(f"  Adam time : {adam_time:.1f}s  "      # Adam timing
      f"({Nepochs_ADAM} epochs)")
if Nbfgs > 0:                                 # BFGS timing (if used)
    print(f"  BFGS time : {bfgs_time:.1f}s  "
          f"({Nbfgs} iters, {optimizer_label})")
print(f"  Total time: "                        # total wall time
      f"{adam_time + bfgs_time:.1f}s")
print(f"Results saved to {RESULTS_DIR}/")      # output directory
\end{lstlisting}


% ======================================================================
\section{Summary of the Training Pipeline}
\label{sec:summary}
% ======================================================================

\begin{center}
\renewcommand{\arraystretch}{1.4}
\begin{tabular}{@{}lcc@{}}
\toprule
& \textbf{Phase 1: Adam} & \textbf{Phase 2: Quasi-Newton} \\
\midrule
Type       & First-order           & Second-order (BFGS / L-BFGS-B) \\
Canonical iters & 2\,000           & 10\,000 \\
Per-step cost & $O(n)$             & $O(n^2)$ (dense) or $O(nm)$ (L-BFGS) \\
Convergence & Linear               & Superlinear \\
Key idea   & Cheap gradient steps   & Curvature-informed steps \\
Hessian    & ---                    & Dense $H_k$ or $m$ correction pairs \\
Scaling    & LR decay               & Self-scaled $\tau_k$ controls $\kappa(H_k)$ \\
Resampling & RAD every 500 epochs   & RAD every \texttt{Nchange} iters \\
\bottomrule
\end{tabular}
\end{center}

The overall flow:
\[
  \underbrace{\text{Adam ($N_A$ iters)}}_{\text{get into a good basin}}
  \;\longrightarrow\;
  \underbrace{\text{SSBroyden2 ($N_B$ iters)}}_{\text{converge to high accuracy}}
\]

where $N_A + N_B = 12{,}000$ (total budget fixed across all experiments).


% ======================================================================
\section{Experiment Suite}
\label{sec:experiments}
% ======================================================================

The YAML configuration system supports 17 experiment configs grouped
into 6 axes.  All share identical fixed parameters (network, sampling,
loss weights, equation, domain, seed) and the same total iteration
budget of 12\,000.

\subsection{Group A: Optimizer Type}

\begin{center}
\begin{tabular}{@{}llp{6.5cm}@{}}
\toprule
\textbf{Config} & \textbf{Method} & \textbf{Purpose} \\
\midrule
\texttt{adam\_only}     & Adam (12\,000 iters) & First-order baseline (floor) \\
\texttt{bfgs\_vanilla}  & Standard BFGS       & No self-scaling baseline \\
\texttt{ssbfgs\_ol}     & SSBFGS (Oren--Luenberger) & OL self-scaling \\
\texttt{ssbfgs\_ab}     & SSBFGS (Al-Baali)    & AB self-scaling \\
\texttt{ssbroyden1}     & SSBroyden $\phi{=}1$ & DFP direction \\
\texttt{canonical}       & SSBroyden2           & Paper default (optimised $\phi$) \\
\texttt{lbfgs}           & L-BFGS-B ($m{=}20$)  & Limited-memory variant \\
\bottomrule
\end{tabular}
\end{center}

\subsection{Group B: Adam Warmup Duration}

\begin{center}
\begin{tabular}{@{}lccp{5cm}@{}}
\toprule
\textbf{Config} & \textbf{Adam} & \textbf{BFGS} & \textbf{Purpose} \\
\midrule
\texttt{warmup\_0}    & 0       & 12\,000 & Cold-start BFGS \\
\texttt{warmup\_500}  & 500     & 11\,500 & Minimal warmup \\
\texttt{canonical}     & 2\,000  & 10\,000 & Paper default \\
\texttt{warmup\_5000} & 5\,000  & 7\,000  & Extended warmup \\
\bottomrule
\end{tabular}
\end{center}

\subsection{Group C: Power Transform}

The power transform minimises $\mathcal{L}^{1/p}$ to improve
Hessian conditioning (Eq.~\eqref{eq:kappa}):

\begin{center}
\begin{tabular}{@{}lcp{5cm}@{}}
\toprule
\textbf{Config} & $p$ & \textbf{Effect} \\
\midrule
\texttt{canonical} & 1 & No transform (raw loss) \\
\texttt{power\_2}  & 2 & $\kappa \to \sqrt{\kappa}$ \\
\texttt{power\_4}  & 4 & $\kappa \to \kappa^{1/4}$ \\
\bottomrule
\end{tabular}
\end{center}

\subsection{Group D: Hessian Initialisation}

\begin{center}
\begin{tabular}{@{}lp{8cm}@{}}
\toprule
\textbf{Config} & \textbf{Effect} \\
\midrule
\texttt{canonical}       & $H_0 = I$ (identity) \\
\texttt{hessian\_scaled} & $H_0 \leftarrow \frac{\bm{y}^\top\bm{s}}{\bm{y}^\top\bm{y}}\,I$ (Nocedal--Wright) \\
\bottomrule
\end{tabular}
\end{center}

\subsection{Group E: Adam Learning Rate}

\begin{center}
\begin{tabular}{@{}lcp{5cm}@{}}
\toprule
\textbf{Config} & $\eta_0$ & \textbf{Purpose} \\
\midrule
\texttt{adam\_lr\_low}  & $10^{-4}$ & Conservative warmup \\
\texttt{canonical}       & $10^{-3}$ & Paper default \\
\texttt{adam\_lr\_high} & $10^{-2}$ & Aggressive warmup \\
\bottomrule
\end{tabular}
\end{center}

\subsection{Group F: BFGS Batch Size (Resampling Frequency)}

\begin{center}
\begin{tabular}{@{}lcp{5cm}@{}}
\toprule
\textbf{Config} & \texttt{Nchange} & \textbf{Purpose} \\
\midrule
\texttt{bfgs\_batch\_200}  & 200   & Frequent RAD, less Hessian maturity \\
\texttt{canonical}          & 500   & Paper default \\
\texttt{bfgs\_batch\_1000} & 1\,000 & Rare RAD, better Hessian convergence \\
\bottomrule
\end{tabular}
\end{center}


% ======================================================================
\section{Key Differences from the Burgers Notebook}
\label{sec:differences}
% ======================================================================

\begin{center}
\renewcommand{\arraystretch}{1.3}
\begin{tabular}{@{}lll@{}}
\toprule
\textbf{Aspect} & \textbf{Burgers} & \textbf{Cahn--Hilliard} \\
\midrule
Spatial dimensions   & 1-D ($x$)              & 2-D ($x, y$) \\
PDE order            & 2nd ($u_t, u_x, u_{xx}$) & 4th ($U_{xxxx}$, etc.) \\
Autograd calls       & 2                       & 8 \\
Input features       & $(t, x)$ = 2            & $(t, x, y)$ = 3 \\
Parameters           & 1\,341                  & 5\,121 \\
Hessian memory       & $\sim$14 MB             & $\sim$210 MB \\
Initial condition    & $-\sin(\pi x)$ (closed form) & Random field (interpolated) \\
Boundary conditions  & Dirichlet ($u=0$)       & Periodic (value + derivative) \\
Configuration        & Hardcoded                & YAML-driven \\
Nonlinearity         & $u\,u_x$ (convection)   & $\nabla^2(U^3)$ (quartic) \\
\bottomrule
\end{tabular}
\end{center}


% ======================================================================
\appendix
\section{Notation Reference}
\label{sec:notation}
% ======================================================================

\begin{center}
\renewcommand{\arraystretch}{1.3}
\begin{tabular}{@{}cl@{}}
\toprule
\textbf{Symbol} & \textbf{Meaning} \\
\midrule
$U(t,x,y)$                    & Scalar order parameter (phase field) \\
$\hat{U}(t,x,y;\bm\theta)$   & Neural network approximation \\
$\bm\theta$                    & All trainable parameters ($\in\mathbb{R}^{5121}$) \\
$D$                            & Diffusion/mobility constant (= 1) \\
$a_2$                          & Quadratic chemical potential coefficient (= 1) \\
$a_4$                          & Quartic chemical potential coefficient (= 1) \\
$\nabla^2$                     & Laplacian: $\partial_{xx} + \partial_{yy}$ \\
$\nabla^4$                     & Biharmonic: $\partial_{xxxx} + 2\partial_{xxyy} + \partial_{yyyy}$ \\
$f(t,x,y;\bm\theta)$          & PDE residual (Eq.~\eqref{eq:residual}) \\
$\mathcal{L}$                 & Total PINN loss (PDE + IC + BC) \\
$\lambda_{\mathrm{pde}},\lambda_{\mathrm{ic}},\lambda_{\mathrm{bc}}$
  & Loss weights (1, 5, 5) \\
$H_k$                          & Inverse Hessian approximation at iteration $k$ \\
$\tau_k$                       & Self-scaling factor \\
$\phi_k$                       & Broyden family interpolation parameter \\
$\bm{s}_k, \bm{y}_k$          & Step and gradient difference vectors \\
$\rho_k$                       & $1 / (\bm{y}_k^\top \bm{s}_k)$ \\
$p$                            & Power transform exponent ($\mathcal{L}_{\mathrm{eff}} = \mathcal{L}^{1/p}$) \\
$\kappa$                       & Condition number of the Hessian \\
$U_0(x,y)$                    & Initial condition (smoothed random field) \\
$G_{\sigma}$                  & Gaussian kernel with std.\ dev.\ $\sigma$ \\
\bottomrule
\end{tabular}
\end{center}


\end{document}
